%% kapa: exam=1 round=26 year=2015 period=2 form=A subject=appraisal_law subject_ko=부동산관계법규 seq=14 number=14
%% source: Q-Net 감정평가사 기출문제 (cst00309) artlSeq=5138001 file=extracted/1008/5138001/제26회_감정평가사 1차_기출/제26회_감정평가사 1차_기출/2교시_감평_A형.pdf
%% conversion: kordoc 4.24.51 -> tools/gen1.py (auto); review=proofread-sonnet; flags=none
\begin{question}{14}
국토의 계획 및 이용에 관한 법령상 도시ㆍ군관리계획에 관한 설명으로 옳지 않은 것은?
\begin{choices}
  \item 이해관계인을 포함한 주민은 기반시설의 정비에 관한 사항에 대하여 도시ㆍ군관리계획의 입안권자에게 도시ㆍ군관리계획의 입안을 제안할 수 있다.
  \item 국토교통부장관은 국가계획과 관련된 경우에는 직접 도시ㆍ군관리계획을 입안할 수 있다.
  \item 도시ㆍ군관리계획입안의 제안을 받은 입안권자는 부득이한 사정이 있는 경우를 제외하고는 제안일부터 30일 이내에 도시ㆍ군관리계획입안에의 반영 여부를 제안자에게 통보하여야 한다.
  \item 개발제한구역의 지정 및 변경에 관한 도시ㆍ군관리계획은 국토교통부장관이 결정한다.
  \item 도시ㆍ군관리계획 결정의 효력은 지형도면을 고시한 날부터 발생한다.
\end{choices}
\end{question}
\answer{3}
